% !TeX root = 20261102_exp93_fullholdout.tex
% !LW recipe=lablog XeLaTeX
\input{_support/style.tex}
\usepackage{tabularx}
\usepackage{tikz}
\usetikzlibrary{arrows.meta,positioning,shapes.geometric,calc}
\newcolumntype{Y}{>{\raggedright\arraybackslash}X}
\setlist[itemize]{leftmargin=*,itemsep=2pt,topsep=3pt}
\hypersetup{pdftitle={EXP93: 기존 CSV와 PCAP--netFound CIC2018의 구성·문제1·문제2 대조}}
\begin{document}
\reporttitle{2026-11-02 실행 · 2026-11-03 재작성}{EXP93: 기존 CSV와 PCAP--netFound CIC2018의\\구성 · 문제1 · 문제2 대조와 재검토}
\recordnote{실행 \path{20261102_003121_cic2018_pcap_expanded_d1f3d5cf}은 120/120개 조건을 반환값 0으로 완료했고 수치는 보존한다. 11월 3일 재검토에서 Bot family 라벨이 Bot 피해 host가 아닌 capture(172.31.69.13)에 alias를 걸어 만들어진 것임을 PCAP 스캔으로 확인했다. Bot 수치, Bot을 포함한 macro-F1, benign/bot 충돌 수는 데이터 성질의 근거로 쓰지 않는다. 본문은 용어 정의 $\to$ 구성 대조 $\to$ 문제1 $\to$ 문제2 $\to$ Bot 결함 $\to$ 재실행 protocol 순이다.}

\section{용어와 번호 규약}
\begin{itemize}
\item \textbf{문제1 = 동일 벡터·상이 라벨.} 분류기 입력이 완전히 같은 행에 서로 다른 family 라벨이 붙은 사례. 결정적 분류기는 둘을 동시에 맞힐 수 없다.
\item \textbf{문제2 = head 조기 포화.} DDoS·BF 같은 head family가 class당 수 개의 라벨로 이미 90\%대 F1에 도달해 추가 라벨의 개선 여지가 작다는 관찰. 2026-10-29 슬라이드는 번호가 반대이므로 인용 시 주의한다.
\item \textbf{PCAP에는 라벨이 없다.} CIC2018은 host별 원시 패킷 캡처와 CICFlowMeter CSV를 배포한다. 라벨은 CSV에만 있고, 그 라벨도 ``시간창 안에서 공격자 IP와 피해자 IP가 통신하면 공격''이라는 규칙으로 만들어졌다. EXP93은 PCAP에서 flow를 재구성한 뒤 CNS2022 교정 규칙을 다시 적용해 라벨을 생성했다.
\item \textbf{라벨 규칙 = 시간창 + endpoint (+ 일부 조건).} 예: SSH brute force는 ``02-14 14:01--15:31, 13.58.98.64 $\to$ 172.31.69.25:22''. 일부 규칙에만 payload·RST 조건이 붙는다.
\item \textbf{증거.} 그 공격을 정의하는 행위가 패킷 안에 관측된 것. brute force는 인증 payload 교환, DoS는 요청 전송, web은 HTTP 요청 내용, 스캔은 반대로 payload 없는 SYN과 무응답/RST. 증거는 공격마다 다르다.
\item \textbf{창 안·증거 없음 (= Attempted, unresolved).} 시간창과 endpoint는 맞지만 증거가 없는 flow. 예: 창 안의 공격자$\to$victim:22 연결이 handshake만 하고 끝남, 포트가 닫혀 SYN$\to$RST만 반복된 FTP brute force, C2가 멈춘 뒤 slave의 SYN$\to$RST 재접속. 패킷만 보면 거절된 정상 접속과 같다.
\item \textbf{라벨 상태 세 가지.} \textbf{confirmed} = 창·endpoint·증거 모두 충족. \textbf{unresolved} = 창·endpoint는 충족하나 증거 없음(현재 corpus는 Attempted 규칙이 있는 경우만 benign으로 매핑, 나머지는 측정 없이 benign). \textbf{background} = 어떤 규칙에도 걸리지 않은 그날 host의 나머지 트래픽, benign으로 라벨.
\item \textbf{group (문제1 감사 단위).} class를 묶은 것이 아니다. 입력이 byte 단위로 완전히 같은 flow들의 집합이다. 예: SYN 1개짜리 flow 300개가 header·timing 토큰까지 모두 같으면 group 1개(행 300개). 그 300개 안에 benign과 inf가 섞여 있으면 ``cross-label group''이다.
\item \textbf{retained.} reservoir sampling으로 corpus에 남긴 flow. PCAP 통계의 ``raw''는 split 직후 전량, ``retained''는 남긴 수다.
\end{itemize}

\begin{figure}[H]
\centering
\resizebox{\linewidth}{!}{%
\begin{tikzpicture}[font=\small, >=Latex, node distance=5mm and 9mm,
  q/.style={diamond, aspect=2.6, draw, inner sep=1pt, align=center, fill=gray!8},
  s/.style={rectangle, rounded corners, draw, align=center, minimum width=26mm, minimum height=7mm},
  conf/.style={s, fill=red!12}, unres/.style={s, fill=yellow!22}, back/.style={s, fill=green!12}, bnd/.style={s, fill=gray!18}]
\node[s, fill=blue!8] (flow) {PCAP에서 재구성한 flow\\(양방향 5-tuple, 시작·종료 시각, payload·flag 합계)};
\node[q, below=of flow] (ep) {endpoint 일치?\\공격자 IP $\leftrightarrow$ 피해자 IP (+port)};
\node[q, below=of ep] (tw) {시간창 안?\\start$\ge$창 시작, end$\le$창 종료};
\node[q, below=of tw] (ev) {증거 조건?\\payload$>$0, RST 없음, 패킷 수 등};
\node[conf, below=of ev] (c) {\textbf{confirmed}\\family 라벨 (학습·평가)};
\node[unres, right=12mm of ev] (u) {\textbf{unresolved} (= Attempted)\\창·endpoint 일치, 증거 없음\\현재: benign 또는 background};
\node[bnd, right=12mm of tw] (b) {\textbf{boundary}\\창 경계에 걸침, 제외};
\node[back, left=10mm of tw] (bg) {\textbf{background}\\나머지 트래픽\\benign 라벨};
\draw[->] (flow) -- (ep);
\draw[->] (ep) -- node[right,xshift=1mm]{예} (tw);
\draw[->] (ep.west) -| node[above,pos=.25]{아니오} (bg);
\draw[->] (tw) -- node[right,xshift=1mm]{예} (ev);
\draw[->] (tw.east) -- node[above]{걸침} (b);
\draw[->] (tw.west) -- node[above]{창 밖} (bg);
\draw[->] (ev) -- node[right,xshift=1mm]{예} (c);
\draw[->] (ev.east) -- node[above]{아니오} (u);
\end{tikzpicture}}
\caption{라벨 규칙의 판정 순서. 색은 라벨 상태(빨강 confirmed, 노랑 unresolved, 초록 background, 회색 boundary). 창 밖 또는 endpoint 불일치는 모두 background benign이 된다.}
\end{figure}

\begin{figure}[H]
\centering
\begin{tikzpicture}[font=\small, >=Latex]
\draw[->, thick] (0,0) -- (13.2,0) node[below left]{시각 (하루 capture)};
\fill[red!10] (4,0) rectangle (9,1.9);
\draw[red!60!black, thick] (4,0) -- (4,1.9); \node[red!60!black, font=\scriptsize, below] at (4,0) {창 시작};
\draw[red!60!black, thick] (9,0) -- (9,1.9); \node[red!60!black, font=\scriptsize, below] at (9,0) {창 종료};
\node[red!60!black] at (6.5,2.3) {공격 규칙 시간창 (공격자 IP $\leftrightarrow$ 피해자 IP)};
\foreach \x/\y/\c/\t in {1.0/1.5/green!60!black/background, 2.2/0.5/green!60!black/background,
                         5.0/1.5/red!70!black/confirmed, 6.2/0.5/orange!90!black/unresolved, 7.4/1.5/red!70!black/confirmed,
                         7.9/0.5/orange!90!black/unresolved, 10.8/1.5/green!60!black/background, 12.0/0.5/green!60!black/background}
  {\draw[\c, line width=2.2pt] (\x,\y) -- ++(0.8,0); \node[\c, font=\scriptsize, above] at (\x+0.4,\y) {\t};}
\draw[gray, line width=2.2pt] (3.6,1.0) -- ++(0.8,0); \node[gray, font=\scriptsize, above] at (4.0,1.0) {boundary};
\draw[green!60!black, line width=2.2pt] (5.6,1.0) -- ++(0.8,0); \node[green!60!black, font=\scriptsize, above] at (6.0,1.0) {background};
\node[font=\scriptsize, align=left, anchor=north west, text width=12.8cm] at (0,-0.55) {가로선 하나 = flow 하나(시작$\to$종료). 창 안이라도 endpoint가 다른 flow(예: 피해 host의 DNS)는 background, 증거가 없는 flow(예: handshake만)는 unresolved.};
\end{tikzpicture}
\caption{같은 host의 하루 트래픽에서 시간창 규칙이 flow를 나누는 방식.}
\end{figure}

\begin{table}[H]
\centering\small
\caption{EXP93 규칙의 실제 예. ``증거 없음'' 열은 창·endpoint는 맞지만 증거 조건을 통과하지 못한 flow가 현재 corpus에서 어디로 가는지다.}
\footnotesize
\begin{tabularx}{\linewidth}{>{\raggedright\arraybackslash}p{.14\linewidth}>{\raggedright\arraybackslash}p{.24\linewidth}>{\raggedright\arraybackslash}p{.18\linewidth}YY}
\toprule 공격 & 시간창 (UTC) · endpoint & 증거 조건 & 증거 없음의 예 & 현재 처리\\\midrule
SSH brute-force & 02-14 14:01--15:31\newline 13.58.98.64 $\to$ .25:22 & 공격자 payload $\ge$1 & handshake만, SYN$\to$RST & Attempted 규칙 $\to$ benign\\
FTP brute-force & 02-14 10:32--12:09\newline 18.221.219.4 $\to$ .25:21 & (규칙 없음) & 포트 닫힘, 전량 SYN$\to$RST/ACK & 규칙 자체가 없어 background benign\\
GoldenEye & 02-15 09:26--10:09\newline 18.219.211.138 $\to$ .25 & 총 payload $\ge$1 그리고 (RST 없음 또는 5.05초 이상) & 요청 전에 RST로 끊긴 연결 & 규칙 불일치 $\to$ background benign\\
Web (XSS) & 02-22 13:50--14:29\newline 18.218.115.60 $\to$ .28 & 지정 source port 제외 & 공격자 브라우저의 정적 파일 요청 & 제외 port만 benign, 나머지는 web\\
NMAP (inf) & 02-28 10:50--14:40\newline .24 $\to$ 내부 21 host & (없음, port 68만 제외) & 창 안 victim의 정상 SMB·DNS & 전부 inf\\
Bot & 03-02 14:13--19:54\newline $\leftrightarrow$ 18.219.211.138 & 총 payload $\ge$1 & 빈 payload, SYN$\to$RST & Attempted $\to$ benign. 단 host가 틀려 전체가 결함(\S6)\\
\bottomrule
\end{tabularx}
\end{table}

\section{구성 대조}
\begin{table}[H]
\centering\small
\caption{train / test별 family 행 수. CSV는 EXP89의 정제 전 \texttt{original} split, PCAP은 EXP93 group split. 행 순서는 PCAP full test 수 내림차순.}
\begin{tabular}{llrrr}
\toprule
구분 & family & CSV (CICFlowMeter) & PCAP raw & PCAP retained\\\midrule
train & ddos   & 794,609 & 784,396 & 1,024\\
      & benign & 10,508,775 & 223,250 & 1,024\\
      & inf    & 112,891 & 132,046 & 1,024\\
      & bf     & 345,116 & 65,874 & 1,024\\
      & dos    & 181,779 & 47,288 & 1,024\\
      & bot    & 124,621 & 13,972 & 1,024\\
      & web    & 1,522 & 893 & 893\\
      & 합계   & 12,069,313 & 1,266,719 & 7,037\\\midrule
test  & ddos   & 264,872 & 167,096 & 167,096\\
      & benign & 3,502,926 & 47,103 & 47,103\\
      & inf    & 37,631 & 28,192 & 28,192\\
      & bf     & 115,040 & 14,379 & 14,379\\
      & dos    & 60,594 & 10,304 & 10,304\\
      & bot    & 41,542 & 2,712 & 2,712\\
      & web    & 509 & 184 & 184\\
      & 합계   & 4,023,114 & 269,970 & 269,970\\
\bottomrule
\end{tabular}
\end{table}

\begin{itemize}
\item \textbf{retained train이 있는 이유.} frozen netFound 인코딩 비용을 평가 집단(test) 밖에서 줄이기 위해 train/validation만 family당 1,024개로 자른다(protocol의 \texttt{holdout\_policy}). few-shot은 이 pool에서 class당 n$\le$50개만 뽑으므로 학습곡선에는 영향이 없다.
\item \textbf{retained train의 부작용.} PCA가 7,057행에만 적합되고, 문제1 감사가 보존 corpus 283,367행에 한정된다. 같은 날 같은 host의 benign--공격 충돌 짝이 표본화로 사라져 충돌률이 과소 집계된다(\S4).
\item \textbf{validation 1,024개는 이번 분류기에서 쓰지 않았다.}
\item \textbf{retained는 무작위 표본이다.} seed 43의 reservoir sampling(Algorithm R)으로 family$\times$split마다 stream 전체에서 균등 확률로 1,024개를 남기고 flow\_id 순으로 저장한다. 시간·host·annotation에 따른 층화는 없다.
\item \textbf{출처 차이의 원인은 benign이다.} 공격은 그날 피해 host 1대에 집중되므로 공격 행 수는 같은 자릿수다(train CSV 1,560,538 vs PCAP 1,044,469, test 520,188 vs 222,867). benign은 CSV가 전체 host(수백 대)·10일, PCAP이 host 1대·7일이라 train 47배, test 74배 차이가 난다. 표1의 10배·20배 격차는 이 benign 차이가 만든 것이다.
\item \textbf{다운로드 방식.} 공식 안내는 \texttt{aws s3 sync}로 bucket \texttt{cse-cic-ids2018} 전체를 받는 것이다. EXP93은 같은 bucket의 공개 HTTPS URL에서 하루 \path{pcap.zip}(40--45 GB, host별 capture 약 440개)의 중앙 디렉터리를 HTTP Range로 읽고 host member 하나만 추출했다. ETag를 If-Match로 고정하고 ZIP CRC32·SHA256을 receipt(\path{fetch.json})에 남겼으므로 그 member는 s3 sync로 받은 파일과 byte 단위로 같다. 다만 부분 수집이다. 예: 02-21 DDoS part1은 18.3 GB 파일을 위해 2.06 GB를 전송했다.
\item \textbf{왜 host 1대·7일인가.} 7개 family를 각각 capture 하나로 덮는 최소 수집이었다. 02-16(dos), 02-20(ddos), 03-01(inf)은 이미 덮인 family의 둘째 날이라 제외됐다. 이 선택은 protocol·commit에 사유로 기록되지 않았고, 기술적 실패가 아니라 범위 선택이다. 같은 방식으로 member를 추가하면 확장할 수 있다.
\end{itemize}

\begin{table}[H]
\centering\small
\caption{CIC2018 원 라벨 15종이 EXP93 corpus에 들어간 경로. Attempted 때문에 빠진 것은 FTP 하나이고, 나머지 누락은 날짜를 내려받지 않은 선택이다. 02-16·02-20·03-01은 이미 덮인 family의 둘째 날이라 수집 범위에서 뺐고, 그 사유는 protocol에 기록되지 않았었다.}
\begin{tabularx}{\linewidth}{l>{\raggedright\arraybackslash}p{.3\linewidth}>{\raggedright\arraybackslash}p{.22\linewidth}Y}
\toprule family & CIC 원 라벨 & EXP93 상태 & 이유\\\midrule
ddos & DDOS-HOIC, DDoS-LOIC-UDP & 포함 (02-21) & \\
     & DDoS-LOIC-HTTP & 제외 & 02-20 capture 미수집\\
dos  & GoldenEye, Slowloris & 포함 (02-15) & \\
     & Hulk, SlowHTTPTest & 제외 & 02-16 capture 미수집\\
bf   & SSH-Bruteforce & 포함 (02-14) & \\
     & FTP-BruteForce & 제외 & CNS2022: 포트 닫힘, 전량 Attempted. EXP93은 규칙을 두지 않아 background benign\\
web  & Brute Force-Web, XSS, SQL Injection & 포함 (02-22, 02-23) & \\
inf  & Infilteration & 02-28 host .24만 포함 & 03-01 host .13 capture 미수집\\
bot  & Bot & 포함, 결함 & 피해 host가 아닌 .13 capture에 alias(\S6)\\
benign & Benign & background로 대체 & 규칙 미매칭 트래픽. CSV Benign과 정의가 다름\\
\bottomrule
\end{tabularx}
\end{table}

\begin{table}[H]
\centering\small
\caption{데이터 구성의 차이.}
\begin{tabularx}{\linewidth}{p{.15\linewidth}YY}
\toprule 항목 & CSV (CICFlowMeter) & PCAP--netFound (EXP93)\\\midrule
출처 & 10일, 전체 host, 16,232,943행 & 7일, 9개 capture, 날짜당 피해 host 1대, 1,812,228 flow\\
flow 정의 & 양방향, active timeout 120초 & 양방향 5-tuple, idle 120초, \textbf{active 5초}(EXP92는 120초, 미기록이었음), RST·양쪽 FIN·새 SYN 종료\\
입력 & 78개 고정 통계(실험은 NF-v3 46개) & netFound-v2-small 512차원 $\to$ training PCA 64차원. 토큰: IP header 길이·TOS·length·TTL, TCP flag·window·상대 seq/ack, burst 시간. payload·주소·port·절대시각 제외\\
라벨 & CIC 원 라벨 15종 $\to$ 7 family & CNS2022 규칙 $\to$ 7 family. 미매칭 = background benign. Attempted = benign\\
family 범위 & DDoS 3, DoS 4, BF 2(FTP 포함), Web 3, Bot, Inf & DDoS 2(LOIC-HTTP 제외), DoS 2(Hulk·SlowHTTPTest 제외), BF는 SSH만, Web 3, Bot(결함), Inf(02-28 host만)\\
분할 & 승인된 시간순 outer split & SHA256(5-tuple group) 70/15/15. 연결마다 group이 달라 같은 분(分)의 flow가 train/test에 섞임\\
\bottomrule
\end{tabularx}
\end{table}

\section{라벨 상태 수와 ``창 안·증거 없음''의 실제 규모}
\begin{itemize}
\item \textbf{corpus label status.} background benign 317,603 · confirmed attack 1,491,948 · attempted-to-benign 2,646 · boundary 19 · incomplete capture 12.
\item \textbf{2,646은 전체 CIC2018이 아니다.} 이 9개 capture에서, 명시적 Attempted 규칙이 있는 두 경우(SSH payload 0, Bot RST·빈 payload)만 센 수다.
\item \textbf{세지 않은 창 안·증거 없음.} FTP brute force(02-14, 포트 닫힘)는 규칙이 없어 전량 background benign. GoldenEye의 RST 조건을 통과하지 못한 flow도 규칙 불일치 뒤 background benign. inf 규칙은 payload 조건이 없어 창 안 모든 내부 통신이 inf.
\item \textbf{따라서} 창 안·증거 없음 flow의 상당수가 현재 benign class에 섞여 있고, 그 수는 측정되지 않았다. 재실행에서는 이 집합을 unresolved로 분리해 family별로 센다(\S8).
\end{itemize}

\section{문제1: 동일 벡터·상이 라벨}
\begin{table}[H]
\centering\small
\caption{family별 cross-label 현황. CSV는 78개 통계 정확 일치(16,232,943행 전체). PCAP은 netFound 입력 signature 일치(보존 corpus 283,367행). 입력 정의·범위·라벨 규칙이 달라 비율은 직접 비교하지 않는다.}
\begin{tabularx}{\linewidth}{lrrrrY}
\toprule
& \multicolumn{3}{c}{CSV} & \multicolumn{2}{c}{PCAP(보존 corpus)}\\
family & 행 & 충돌 행 & 비율(\%) & cross-label group & 상대 family(group)\\\midrule
benign & 13,484,708 & 125,771 & 0.93 & 49 & bot 43 · inf 6\\
ddos & 1,263,933 & 0 & 0.00 & 0 & --\\
dos & 654,300 & 77,277 & 11.81 & 0 & --\\
bf & 380,949 & 162,479 & 42.65 & 0 & --\\
bot & 286,191 & 0 & 0.00 & 43 & benign\\
inf & 161,934 & 41,279 & 25.49 & 6 & benign\\
web & 928 & 5 & 0.54 & 0 & --\\\midrule
합계 & 16,232,943 & 406,811 & 2.51 & 49 & 507행, 0.179\%\\
\bottomrule
\end{tabularx}
\end{table}

\textbf{49 group과 431 group.}
\begin{itemize}
\item \textbf{49 group} = netFound \emph{입력 signature}(flow\_duration, burst 토큰, 방향, 바이트, iat, 패킷 수, protocol)를 hash했을 때 완전히 같은 flow 집합 중, 서로 다른 family가 섞인 집합의 수. 해당 행은 507개.
\item \textbf{431 group} = \emph{PCA 64차원 vector}가 byte 단위로 같은 flow 집합(18,454개) 중, family가 섞인 집합의 수. 행 수는 예시 100 group까지만 저장해 알 수 없다.
\item \textbf{왜 431 $>$ 49인가.} signature의 flow\_duration은 encoder에 전달되지 않고, iat는 tokenizer에서 ms로 절삭된다. signature가 다른 flow가 같은 모델 입력으로 수렴하므로, 모델 기준 충돌은 431이 맞고 49는 하한이다.
\item \textbf{보존 corpus 한정.} train/validation을 family당 1,024개로 자른 뒤 감사했다. 같은 날 같은 host의 benign--공격 짝 대부분이 사라졌으므로 0.179\%는 과소 집계다. 전체 181만 flow의 hash 감사는 streaming으로 가능하다.
\item \textbf{충돌 구조가 다르다.} CSV는 bf·inf·dos에서 크고 bot은 0. PCAP은 bot·inf에서만 남았다. CSV의 bf·dos 충돌은 flow 종료 처리·특징 계산 결함(Engelen 2021, Rosay 2022)이며 새 pipeline이 고쳤다. PCAP의 benign/bot은 \S6의 라벨 결함, benign/inf는 NMAP 창 안팎의 같은 내부 통신이다.
\end{itemize}

\begin{table}[H]
\centering\small
\caption{충돌이 생기는 원인 다섯 가지. ①②는 라벨 규칙 오류(줄일 수 있음), ③은 원리적 한계(분리·정량화), ④⑤는 pipeline·encoder 선택.}
\begin{tabularx}{\linewidth}{p{.03\linewidth}p{.15\linewidth}YYp{.22\linewidth}}
\toprule & 원인 & 예 & EXP93 상태 & 대응\\\midrule
① & 창 과포함 & 창 안 피해 host의 DNS·keepalive·운영자 RDP(131.202.242.193)가 공격으로 라벨 & Bot capture 전체가 이 유형. inf는 payload 조건 없음 & 공격별 증거 조건(port·payload·응답) 추가\\
② & 창 미포함 & C2가 멈춘 뒤 slave의 재접속, 창 시작 전 hello & 미측정 & 창 밖 공격자 endpoint 통신을 unresolved로 분리\\
③ & 관측에 의도가 없음 & SYN 1개 스캔 vs 실패한 정상 접속, payload 없는 brute force 시도 & Attempted 2,646을 benign으로 매핑. FTP·GoldenEye 탈락분은 미측정 & 원리적 한계. unresolved로 제외하고 수·비율 보고. 원시 패킷 동일 집단의 oracle 상한 병기\\
④ & flow 분할 산물 & 순수 ACK 꼬리, 5초 active timeout이 만든 SYN 없는 중간 조각 & ACK 꼬리는 제외. active 5초는 적용 중이며 미기록 & active 120초 복원, end\_reason 분포 보고\\
⑤ & 표현 붕괴 & 서로 다른 원시 flow가 같은 입력으로 수렴(payload 제거, iat ms 절삭, 6패킷$\times$12버스트 상한) & 49 $\to$ 431 group의 격차 & 원시 패킷 동일 $\le$ 토큰 동일 $\le$ 모델 입력 동일의 3층 감사. 증가분이 크면 payload-aware encoder 정당화\\
\bottomrule
\end{tabularx}
\end{table}

\section{문제2: head 조기 포화}
\begin{table}[H]
\centering\small
\caption{정상 포함 n=20 class F1 (\%, 5 seed 평균). CSV는 EXP89 정제 전 \texttt{original} train/test, 원래 전체 test. 행 순서는 PCAP full test 수 내림차순. 각 데이터 안에서 두 모델 중 높은 값을 굵게 표시.}
\begin{tabular}{lrrrrrr}
\toprule
& \multicolumn{3}{c}{CSV (EXP89 original)} & \multicolumn{3}{c}{PCAP (EXP93 full test)}\\
family & test & XGBoost & TabPFN & test & XGBoost & TabPFN\\\midrule
ddos   & 264,872 & \textbf{98.16} & 94.99 & 167,096 & 97.80 & \textbf{98.91}\\
benign & 3,502,926 & 69.73 & \textbf{72.32} & 47,103 & \textbf{81.81} & 81.52\\
inf    & 37,631 & 3.71 & \textbf{3.78} & 28,192 & \textbf{95.32} & 94.47\\
bf     & 115,040 & 90.78 & \textbf{91.54} & 14,379 & 98.29 & \textbf{99.54}\\
dos    & 60,594 & 53.75 & \textbf{64.18} & 10,304 & 83.65 & \textbf{85.19}\\
bot    & 41,542 & 57.77 & \textbf{83.82} & 2,712 & 33.92 & \textbf{39.93}\\
web    & 509 & \textbf{0.42} & 0.39 & 184 & 16.53 & \textbf{52.13}\\\midrule
macro & 4,023,114 & 53.47 & \textbf{58.72} & 269,970 & 72.47 & \textbf{78.81}\\
\bottomrule
\end{tabular}
\end{table}

\begin{itemize}
\item \textbf{head는 두 데이터에서 같다.} DDoS·BF는 n=20에서 양쪽 모두 90\%대 후반. 조기 포화는 특징 추출 방식이 아니라 반복적·대량 공격의 성질이다.
\item \textbf{macro 차이는 tail에서 온다.} PCAP macro-F1이 n=20에 80\% 미만인 것은 Bot(라벨 결함)과 Web(test 184행) 때문. ``조기 포화 조건 미충족''은 이 두 class에 의존한다.
\item \textbf{inf 3.7\% $\to$ 95\%는 과제가 바뀐 것이다.} PCAP inf train의 대다수가 NMAP 스캔 SYN flow라 포트스캔 탐지 성능이다.
\item \textbf{개선 여지의 위치.} head를 어렵게 만드는 것이 아니라 tail(bot·web)과 benign 오탐. 측정 단위는 class별 학습곡선과 데이터 상한선(\S8).
\end{itemize}

\begin{table}[H]
\centering\small
\caption{PCAP full test Macro-F1 (\%, 5 seed 평균 $\pm$ 표본 표준편차). $\Delta$ = TabPFN$-$XGBoost. Bot 포함 값이므로 실행 기록이다.}
\begin{tabular}{r|rrr|rrr}
\toprule
& \multicolumn{3}{c|}{정상 포함} & \multicolumn{3}{c}{공격만}\\
n & XGBoost & TabPFN & $\Delta$ & XGBoost & TabPFN & $\Delta$\\\midrule
2  & 38.66$\pm$3.62 & 55.88$\pm$5.87 & +17.21 & 36.56$\pm$8.73 & 58.46$\pm$7.08 & +21.90\\
3  & 54.97$\pm$2.58 & 65.28$\pm$5.78 & +10.31 & 58.58$\pm$3.80 & 71.38$\pm$7.77 & +12.80\\
5  & 62.80$\pm$4.03 & 72.25$\pm$4.22 & +9.46  & 66.48$\pm$6.83 & 75.36$\pm$4.22 & +8.89\\
10 & 67.45$\pm$3.11 & 77.11$\pm$2.61 & +9.66  & 72.71$\pm$8.00 & 80.92$\pm$4.08 & +8.21\\
20 & 72.47$\pm$3.69 & 78.81$\pm$1.79 & +6.34  & 84.38$\pm$11.00 & 90.27$\pm$5.95 & +5.89\\
50 & 78.63$\pm$3.03 & 80.87$\pm$1.28 & +2.24  & 92.29$\pm$2.72 & 92.62$\pm$2.03 & +0.32\\
\bottomrule
\end{tabular}
\end{table}

\reportfigure{figures/exp93_fulltest_macro_f1.pdf}{EXP93 full test learning curve. 점은 5 seed 평균, 오차막대는 표본 표준편차. 왼쪽 정상 포함, 오른쪽 공격 family만.}
\reportfigure{figures/exp93_fulltest_class_f1_n20.pdf}{n=20 family별 F1. Bot은 \S6의 라벨 결함을 반영한다.}
\reportfigure{figures/exp93_fulltest_vs_exp89.pdf}{Solid = PCAP--netFound full test, dashed = EXP89 정제 전 CSV. 절대 F1을 순위화하지 않고 기울기와 모델 차이를 본다.}

\section{Bot 라벨 결함}
\begin{itemize}
\item \textbf{공식 피해 host.} CIC 일정표의 2018-03-02 Bot 피해 host는 172.31.69.23·17·14·12·10·8·6·26·29·30. EXP93이 고른 \path{capEC2AMAZ-O4EL3NG-172.31.69.13}은 목록에 없다. .13은 전날(03-01) Infiltration 피해 host다.
\item \textbf{alias 도입 경위.} alias 전 실행 \path{20261029_162556_cic2018_pcap_expanded_6f909ba0}의 bot label은 0건. commit \path{0ccd2290}이 ``공개 endpoint 18.219.211.138이 VPC 내부 주소 172.31.69.13으로 관측된다''는 가설로 alias를 추가해 coverage 검사를 통과시켰다.
\item \textbf{alias의 효과.} 172.31.69.13은 capture host 자신이므로 규칙은 ``이 host의 5.7시간 창 안 payload flow 전부 = bot''이 됐다.
\item \textbf{스캔 결과(11월 3일, \path{scripts/v7/exp93/audit_bot_capture.py}).} 아래 표. 피해 host .26에서는 C2가 공개 IP 그대로 port 8080에 15만 패킷으로 보인다. .13에는 C2 패킷이 없고 창 안 트래픽은 CDN HTTPS·SMB·DNS·RDP 배경이다.
\item \textbf{결과에 미친 영향.} bot 19,310 flow는 Windows 배경 트래픽. benign/bot 충돌 43 group, Bot F1 34--40\%, benign F1 81\%는 모두 이 결함의 산물이다.
\end{itemize}

\begin{table}[H]
\centering\small
\caption{Bot 날 host capture 스캔. 창 = CNS2022 Bot 규칙 14:13:28--19:54:52 UTC. C2 = 18.219.211.138, botnet 통신 = port 8080(CNS2022 기록).}
\begin{tabularx}{\linewidth}{p{.27\linewidth}rrrY}
\toprule
capture host & 패킷 & C2 주소 패킷 & port 8080 패킷 & 창 안 상위 port(패킷)\\\midrule
172.31.69.13 (EXP93 사용) & 245,109 & 0 & 8 & 443 (83,408) · 445 (18,305) · 53 (12,728) · 80 (12,413) · 3389 (2,531)\\
172.31.69.26 (공식 피해 host) & 228,488 & 150,393 & 150,400 & 8080 (150,397) · 3389 (61,530) · 445 (3,559)\\
\bottomrule
\end{tabularx}
\end{table}

\section{그 외 재검토 사항}
\begin{itemize}
\item \textbf{분할.} 5-tuple group hash는 같은 공격 세션의 연결들을 train/test에 무작위로 나눈다. 보존 corpus에서 PCA 동일 vector group 18,454개. train$\to$test 정확 중복률 미측정. 프로젝트 규칙(시간순 분할, s05/s06 교훈)과 충돌.
\item \textbf{header 지문.} 토큰에 TTL·TOS·IP length·TCP window 포함. 공격자는 모두 별도 Kali host라 OS·host 지문만으로 갈릴 수 있다. UNSW TTL 지름길과 같은 구조.
\item \textbf{host 1대/일.} benign은 그 host의 미매칭 트래픽. family 범위도 CSV와 다름(\S2).
\item \textbf{active timeout 5초.} 5초 초과 연결이 모두 조각난다. Slowloris·Dropbox·장기 benign 세션이 SYN 없는 중간 조각으로 생성. CICFlowMeter(120초)와 flow 단위가 다름.
\item \textbf{baseline 설정.} XGBoost 300 trees·depth 8을 n=2에도 그대로 적용. n=2--5의 TabPFN 이득 일부는 이 고정 설정의 효과일 수 있다.
\end{itemize}

\section{결론과 재실행 protocol}
\begin{itemize}
\item \textbf{판정.} pipeline·평가 배선은 검증됐으나 두 문제의 근거로는 쓸 수 없다. 문제1은 라벨 결함이 만든 충돌과 표현이 지운 충돌을 분리하지 못했고, 문제2는 head 포화가 그대로 재현됐으며 macro 차이는 결함 class에서 왔다. ``새 표현이 두 문제를 완화했다''는 주장은 보류한다.
\end{itemize}
\textbf{재실행 순서.}
\begin{enumerate}[leftmargin=*,itemsep=2pt]
\item Bot capture를 공식 피해 host(.26 또는 .23)로 교체, alias 삭제, bot 규칙에 port 8080 조건 추가. corpus부터 재구성.
\item 라벨 상태를 confirmed / unresolved / background로 두고 family별 수 보고. FTP·GoldenEye 탈락분도 unresolved.
\item 감사 3층(원시 패킷 동일 $\le$ tokenizer 출력 동일 $\le$ PCA 동일)을 전체 181만 flow에 저장. test의 동일 입력 group에서 다수 라벨을 맞히는 oracle macro-F1을 데이터 상한선으로 병기.
\item 같은 9개 PCAP·같은 flow ID·같은 라벨 위에서 CICFlowMeter 통계 vs netFound 표현의 2$\times$2 비교(EXP92 계획) 복원.
\item 분할은 capture 내부 시간 순방향(앞 70\% train) 기본. web(02-22$\to$02-23)·inf(02-28 host .24 $\to$ 03-01 host .13)는 day-held-out 추가. train$\to$test 정확 중복률 저장.
\item 첫 패킷의 TTL·window·TOS·hl만 쓰는 소형 대조 모델로 header 지문 지름길 측정.
\item active timeout 120초 복원, end\_reason 분포 기록.
\item 1--7 뒤에 ET-BERT·NetMamba 등 payload-aware view를 같은 flow ID·분할 위에 추가.
\end{enumerate}

\sourcepath{tabpfn/results/v7/exp93/20261102\_003121\_cic2018\_pcap\_expanded\_d1f3d5cf/RUN\_META.json}
\sourcepath{tabpfn/results/v7/exp93/20261102\_003121\_cic2018\_pcap\_expanded\_d1f3d5cf/corpus/manifest.json}
\sourcepath{tabpfn/results/v7/exp93/20261102\_003121\_cic2018\_pcap\_expanded\_d1f3d5cf/encoders/netfound-v2-small/audit.json}
\sourcepath{tabpfn/results/v7/exp93/20261029\_162556\_cic2018\_pcap\_expanded\_6f909ba0/corpus/manifest.json}
\sourcepath{docs/lablog/evidence/20261102/exp93\_fulltest\_summary.json}
\sourcepath{docs/lablog/evidence/20261103/exp93\_bot\_capture\_audit.json}
\sourcepath{docs/lablog/evidence/20261030/original\_cic2018\_csv\_audit.json}
\sourcepath{tabpfn/results/v7/exp89/20261027\_144031\_curation\_benign\_n2\_3\_5\_10\_20\_five\_seeds\_e1167f0a/analysis/per\_class.csv}
\sourcepath{docs/lablog/reports/20261027\_exp89.tex}
\sourcepath{configs/experiments/exp93\_protocol.json}
\sourcepath{https://www.unb.ca/cic/datasets/ids-2018.html}
\end{document}
